%=====================================================================
\chapter{Threats to Validity and How to Report Them}\label{ch:validity}
%=====================================================================
\locator{4}{Part~4 --- Measurement, Data and Analysis}

\begin{outcomes}
By the end of this chapter you will be able to:
\begin{tight}
  \item \textbf{Classify} a threat into one of the four validity types.
  \item \textbf{Identify} the threats characteristic of each research design.
  \item \textbf{Distinguish} a threat you mitigated, one you measured, and one
    you merely acknowledged.
  \item \textbf{Write} a threats section that increases rather than decreases a
    reader's confidence.
  \item \textbf{Recognise} the trade-off between validity types, and defend
    your position on it.
\end{tight}
\end{outcomes}

\begin{prereq}
Everything in Parts~III and IV. This chapter is where the designs and the
analysis are audited, and it is the chapter to reread before writing up.
\end{prereq}

\begin{keyterms}
internal validity $\bullet$ external validity $\bullet$ construct validity
$\bullet$ conclusion validity $\bullet$ confounding $\bullet$ selection
$\bullet$ maturation $\bullet$ instrumentation $\bullet$ mono-operation bias
$\bullet$ hypothesis guessing $\bullet$ ecological validity $\bullet$
mitigation
\end{keyterms}

\section{Four Questions About One Study}

Cook and Campbell's four-way taxonomy~\cite{cook1979quasi} remains the standard,
and each type asks a different question about the chain from your treatment to
your conclusion.

\begin{center}
\small
\begin{tabular}{@{}L{3.0cm}L{5.0cm}L{6.3cm}@{}}
\toprule
\textbf{Type} & \textbf{Asks} & \textbf{If it fails} \\
\midrule
Conclusion validity
  & Is there a relationship in the data at all?
  & You claim an effect that is noise, or miss one that is there \\
Internal validity
  & Is the relationship causal?
  & Something other than your treatment produced the effect \\
Construct validity
  & Do the operationalisations capture the intended constructs?
  & You measured or manipulated something other than what you claim \\
External validity
  & Does it hold beyond this setting?
  & The result is true here and nowhere anyone cares about \\
\bottomrule
\end{tabular}
\end{center}

\begin{conceptbox}[Read them in order, because they depend on each other]
The order is not arbitrary. If conclusion validity fails there is no
relationship to be causal about. If internal validity fails, generalising a
non-causal relationship is pointless. If construct validity fails, you do not
know what you generalised.

\medskip
A threats section that opens with external validity --- ``our sample was only
from one company'' --- and never addresses whether the effect exists or is
causal has skipped the two questions a reviewer asks first.
\end{conceptbox}

\section{The Threats, by Type}

\begin{center}
\small
\begin{longtable}{@{}L{3.2cm}L{4.6cm}L{6.4cm}@{}}
\toprule
\textbf{Threat} & \textbf{What it is} & \textbf{Typical mitigation} \\
\midrule
\endfirsthead
\multicolumn{3}{@{}l}{\emph{continued}}\\
\toprule
\textbf{Threat} & \textbf{What it is} & \textbf{Typical mitigation} \\
\midrule
\endhead
\bottomrule
\endfoot
\multicolumn{3}{@{}l}{\textbf{Conclusion validity}}\\[2pt]
Low power
  & Too small to detect the effect of interest
  & A-priori power analysis (\chref{sampling}); within-subject design \\
Violated assumptions
  & The test's assumptions do not hold
  & Check residuals; use robust or rank-based methods (\chref{inference}) \\
Fishing
  & Many comparisons, one reported
  & Pre-specify; correct; label exploratory analyses as such \\
Unreliable measures
  & Noise attenuates the relationship
  & Establish reliability; use multi-item scales (\chref{measurement}) \\
Random irrelevancies
  & Uncontrolled setting variation adds noise
  & Standardise conditions; block on known sources \\[4pt]
\multicolumn{3}{@{}l}{\textbf{Internal validity}}\\[2pt]
Selection
  & Groups differed before treatment
  & Randomise; otherwise an identification strategy (\chref{quasiexp}) \\
History
  & Something else happened during the study
  & Control group experiencing the same period \\
Maturation
  & Subjects changed anyway --- learning, fatigue
  & Control group; counterbalance order \\
Instrumentation
  & The measuring process changed mid-study
  & Fix the instrument; if human coders, check drift over time \\
Mortality
  & Drop-out differs between conditions
  & Report attrition per condition; intention-to-treat analysis \\
Diffusion
  & The control group learns the treatment
  & Separate groups; ask afterwards whether it happened \\[4pt]
\multicolumn{3}{@{}l}{\textbf{Construct validity}}\\[2pt]
Inadequate definition
  & The construct was never pinned down
  & Define before operationalising (\chref{questions}) \\
Mono-operation bias
  & One task, one system, one instantiation of the treatment
  & Multiple tasks or systems; report per-object results \\
Mono-method bias
  & One kind of measure for everything
  & Triangulate methods (\chref{mixed}) \\
Hypothesis guessing
  & Subjects infer the expected result and comply
  & Blind subjects to condition and to hypothesis \\
Experimenter expectancy
  & The researcher's expectation influences conduct or coding
  & Blind coding; independent second coder \\
Evaluation apprehension
  & Subjects perform because they are watched
  & Anonymity; assurance that individuals are not assessed \\[4pt]
\multicolumn{3}{@{}l}{\textbf{External validity}}\\[2pt]
Population
  & Subjects unlike the target population
  & Characterise subjects fully; state the target; argue the specific task
    (\chref{experiments}) \\
Setting
  & Laboratory conditions unlike practice
  & Realistic tasks and time pressure; a field study alongside \\
Time
  & The result is period-specific
  & State the technological context; expect decay \\
Object
  & The systems studied are unrepresentative
  & Select objects on stated criteria, and report which \\
\end{longtable}
\end{center}

\section{The Trade-off Is Real and Must Be Owned}

Increasing one kind of validity often decreases another, and no design escapes
this.

\begin{center}
\small
\begin{tabular}{@{}L{3.6cm}L{5.2cm}L{5.3cm}@{}}
\toprule
\textbf{Move} & \textbf{Gains} & \textbf{Costs} \\
\midrule
Tighter laboratory control
  & Internal, conclusion
  & External --- setting realism \\
Field study in one company
  & External --- setting
  & Internal --- no randomisation \\
More measures per construct
  & Construct
  & Conclusion --- more comparisons; and participant fatigue \\
Larger, more diverse sample
  & External --- population; conclusion --- power
  & Internal --- more heterogeneity to control; and cost \\
\bottomrule
\end{tabular}
\end{center}

\begin{conceptbox}[State your position deliberately]
There is no design that maximises all four, so a competent study chooses. What
distinguishes a good threats section is that the choice is stated as a choice.

\medskip
``We prioritised internal validity: a controlled experiment with randomised
assignment on a standardised task. This limits external validity, since the
task is smaller than industrial code and the subjects are graduate students. We
regard this as the right trade at this stage, because no prior work has
established that the effect exists at all; establishing it under controlled
conditions is a precondition for studying it in the field, which we propose as
the natural follow-up.''

\medskip
That paragraph tells an examiner you understood the trade-off rather than fell
into it.
\end{conceptbox}

\section{Writing the Section}

\begin{alertbox}[The three grades of threat, and why the distinction matters]
Not all threats are handled equally, and conflating them is the fastest way to
lose a reader's trust.

\begin{tightnum}
  \item \textbf{Mitigated.} You designed it out. ``Assignment was randomised,
    removing selection.'' Strongest.
  \item \textbf{Measured.} You could not remove it but you quantified it.
    ``Attrition was 8\% and 9\% across conditions, and drop-outs did not differ
    on baseline experience.'' Nearly as strong, and often more honest.
  \item \textbf{Acknowledged.} You can do neither. ``Subjects were graduate
    students; professionals may show a smaller effect, since the checklist
    likely substitutes for experience they already have.'' Acceptable
    \emph{when it states a direction}.
\end{tightnum}

\medskip
A section made entirely of grade 3 reads as a list of excuses. A section that
shows several grade 1s and 2s, and reserves grade 3 for genuine limits, reads
as competence --- and it is the difference \reg{PhD Regs 2024}{\S14.2(iv)} is
reaching for when it asks that shortcomings ``have been addressed''.
\end{alertbox}

\begin{center}
\small
\begin{tabular}{@{}L{3.6cm}L{10.7cm}@{}}
\toprule
\textbf{Do} & \textbf{Do not} \\
\midrule
Order by validity type, conclusion first
  & Open with ``our sample was small'' and stop \\
Give each threat a direction
  & ``This may have affected the results'' --- which way, and by how much? \\
Name the most serious threat explicitly and treat it first
  & Bury it in the middle, where a reviewer will find it and note that you did
    not \\
Distinguish mitigated, measured and acknowledged
  & Present everything as an equal limitation \\
Say what would resolve it
  & Leave the reader to design the follow-up study \\
\bottomrule
\end{tabular}
\end{center}

\begin{worked}[Two versions of the same paragraph]
\textbf{Weak.} ``There are several threats to validity in this study. The
sample size was small, which may have affected the results. The participants
were students, which may limit generalisability. The study was conducted in one
university, which is another limitation. Future work should address these
issues.''

\medskip
Four sentences, no directions, no mitigations, no follow-up. It reads as though
the section were required rather than meant, and it invites every one of the
questions it avoids.

\medskip
\textbf{Strong.} ``\emph{Conclusion validity.} With 34 participants in a
within-subjects design the study had 0.81 power to detect $d = 0.35$, our
pre-registered smallest effect of interest; smaller effects would likely have
been missed. Residuals were right-skewed, so we report Cliff's $\delta$ with
bootstrap intervals alongside the parametric test; both agree.

\medskip
\emph{Internal validity.} Treatment order was counterbalanced by Latin square,
removing order and material effects. Two participants withdrew, one from each
condition, so attrition is unlikely to be differential.

\medskip
\emph{Construct validity.} Defects found is a partial measure of review
quality: it ignores false positives raised, which we did not record. If the
checklist increases false positives, our measure overstates its benefit. This
is the study's most serious limitation and we recommend recording both in
replications.

\medskip
\emph{External validity.} Participants were MS students with a median of 14
months of industrial experience --- closer to junior professionals than to
senior reviewers. We expect the checklist effect to be \emph{smaller} among
senior reviewers, who plausibly already apply an internalised equivalent, so
the reported effect should be read as an upper bound for experienced
populations. The code artefacts were 200--400 line modules; behaviour on large
legacy systems is untested.''

\medskip
Same study. The second is believable, and it names its own worst problem before
a reviewer can.
\end{worked}

\begin{conceptbox}[Naming your own worst threat is a strength]
Scholars fear that identifying a serious limitation invites rejection. The
opposite is closer to true. A reviewer will find it regardless; the only
question is whether you found it first.

\medskip
An author who names it, sizes it, and says what would resolve it demonstrates
command of their own work. An author who omits it is either unaware --- which
is worse --- or hoping the reviewer will be. \reg{PhD Regs 2024}{\S14.3(D)}
puts this plainly: ``the assessment of the results performed by the author must
not be superficial and lacking substance''.
\end{conceptbox}

\section{Threats by Design}

\begin{center}
\small
\begin{tabular}{@{}L{3.2cm}L{5.3cm}L{5.6cm}@{}}
\toprule
\textbf{Design} & \textbf{Strongest} & \textbf{Watch most carefully} \\
\midrule
Controlled experiment
  & Internal
  & External --- setting and population; mono-operation bias \\
Quasi-experiment
  & External --- realistic setting
  & Internal --- the identifying assumption (\chref{quasiexp}) \\
Repository mining
  & Conclusion --- large samples
  & Construct --- what the trace actually measures; external --- open source
    only \\
Survey
  & External --- if sampling permits
  & Construct --- self-report; internal --- no causal warrant \\
Case study
  & Construct --- rich, contextual
  & External --- analytic not statistical (\chref{casestudy}) \\
Design science
  & Relevance
  & Conclusion --- demonstration mistaken for evaluation
    (\chref{designscience}) \\
Benchmarking
  & Conclusion --- if repeated properly
  & External --- workload representativeness (\chref{benchmarking}) \\
ML evaluation
  & Conclusion --- large test sets
  & Internal --- leakage; construct --- label validity (\chref{mleval}) \\
\bottomrule
\end{tabular}
\end{center}

\begin{templatebox}[C11 --- Threats to validity table]
Fill one row per threat. Anything you cannot grade is a threat you have not
thought about yet.

\medskip
\begin{tabular}{@{}L{2.2cm}L{3.6cm}L{2.4cm}L{4.8cm}@{}}
\textbf{Type} & \textbf{Threat} & \textbf{Grade} & \textbf{Direction and response} \\
Conclusion & \fillline{3.4cm} & \fillline{2.2cm} & \fillline{4.6cm} \\[7pt]
Internal   & \fillline{3.4cm} & \fillline{2.2cm} & \fillline{4.6cm} \\[7pt]
Construct  & \fillline{3.4cm} & \fillline{2.2cm} & \fillline{4.6cm} \\[7pt]
External   & \fillline{3.4cm} & \fillline{2.2cm} & \fillline{4.6cm} \\[7pt]
\end{tabular}

\medskip
Grade: M = mitigated by design, Q = quantified, A = acknowledged.\\
Direction: which way would this bias the result, and by roughly how much?\\
\textbf{Most serious threat overall:} \fillline{9cm}
\end{templatebox}

\begin{traditions}{threats to validity}
\tradEMP{The four-way taxonomy, reported by type, with each threat graded and
given a direction.}
\tradDSR{Reframed as evaluation adequacy: does the evaluation support the
claimed contribution, and what would falsify the design knowledge claim?}
\tradQUAL{Replaced by trustworthiness --- credibility, transferability,
dependability, confirmability (\chref{qda}). Do not force the four-way
taxonomy onto qualitative work.}
\tradFORM{Validity is soundness of proof plus realism of assumptions. The
second is where formal results usually fail to transfer, and it deserves a
paragraph.}
\end{traditions}

\begin{lab}[Grade your own threats]
For your own study design:

\begin{tightnum}
  \item Complete template C11 with at least two threats per validity type.
  \item Grade each M, Q or A.
  \item For every A, write the direction of bias in one sentence. If you cannot
    state a direction, you do not yet understand the threat.
  \item Identify the single most serious threat and write the paragraph that
    treats it first.
  \item Count your Ms. If there are none, your design has not been engineered
    --- go back to \chref{experiments} or \chref{quasiexp} and add a control.
\end{tightnum}
\end{lab}

\begin{checkpoint}
\begin{tightnum}
  \item Why should conclusion validity be addressed before external validity?
  \item Distinguish mitigated, measured and acknowledged threats with an
    example of each from your own work.
  \item Why is ``this may have affected the results'' worthless, and what
    replaces it?
  \item Give a design change that would improve external validity while
    damaging internal validity, and say when that trade is right.
\end{tightnum}
\end{checkpoint}

\begin{chaptersummary}
\begin{tight}
  \item Four validity types, addressed in order: conclusion, internal,
    construct, external. They depend on each other.
  \item Each design has characteristic threats; know your design's and treat
    them first.
  \item Validity types trade against each other. Choose deliberately and say
    that you chose.
  \item Grade every threat as mitigated, measured or acknowledged. A section
    that is all acknowledgement reads as excuses.
  \item Every acknowledged threat needs a direction, and ideally a magnitude.
  \item Name your own worst threat before a reviewer does. \S14.3(D) requires
    that the self-assessment not be superficial.
  \item Do not apply this taxonomy to qualitative work; use trustworthiness.
\end{tight}
\end{chaptersummary}

\begin{reviewq}
  \item The four-way taxonomy was developed for field experiments in social
    science. Which computing designs does it fit poorly, and what would a
    taxonomy built for our field look like?
  \item Threats sections are near-universal and near-universally
    formulaic. What incentive produces the formulaic version, and what
    reviewing practice would change it?
  \item Mono-operation bias --- one task, one system --- affects most software
    engineering experiments, because using several is expensive. Is the field's
    tolerance of it justified? Propose a practical standard.
  \item \reg{PhD Regs 2024}{\S14.3(B)} asks about ecological and population
    external validity separately. Give a study strong in one and weak in the
    other, and say which matters more for a dissertation claiming practical
    relevance under \S14.3(D).
\end{reviewq}
